\section{Harness Details}
\label{app:harness}

\subsection{Agent loop}
\label{app:tool-loop}

Each episode runs one paper with one model, opening with a fixed system
message and a task prompt rendered from the paper's dataset row
(Appendix~\ref{app:prompts}), and every request includes the tools of
Appendix~\ref{app:toolset} and an ephemeral status line giving the rounds used and
remaining. A turn may issue several tool calls, which the harness executes in
order, appending one truncated JSON result each. Near the served context
window a compaction step elides older tool results, leaving a placeholder with
the character count and the on-disk path of the elided text. An episode ends
when the model returns no tool call, reaches the round cap, exhausts its
compute budget, overflows the context beyond what elision relieves, or hits an
unrecoverable error, and the harness then issues one final tool-free request
constrained to the report schema of Appendix~\ref{app:report-format}.
Table~\ref{tab:knobs} lists the code defaults, and the launch scripts of
Appendix~\ref{app:models} record the values a sweep set.

\begin{table}[ht]
\centering
\footnotesize
\renewcommand{\arraystretch}{1.15}
\begin{tabular}{@{}>{\raggedright\arraybackslash}p{4.4cm}>{\raggedright\arraybackslash}p{6.2cm}l@{}}
\toprule
Parameter & Value & Where enforced \\
\midrule
Rounds per episode & 300 & agent loop \\
Model output per turn & 32{,}768 tokens & agent loop \\
Reasoning kept from a turn cut at the output cap & 2{,}000 characters & agent loop \\
Length nudges before the final turn & 2 & agent loop \\
Tool result returned to the model & 40{,}000 serialized characters, at least
200 of each string & tool wrapper \\
Retries of a transient tool error & 1 & tool wrapper \\
Compaction trigger & 0.70 of the served context window less the completion
reservation & compaction \\
Recent turns kept verbatim & 20{,}000 tokens, sized at three characters per
token & compaction \\
Tool result elided in the rewritten span & longer than 1{,}000 characters &
compaction \\
Compute ceiling per episode & 8, 32, or 96 H100-hours, the upper edge of the
paper's band (Section~\ref{subsec:construction}) & budget \\
GPUs per session & default 1, limited to the node size, four on the GH200
nodes we use & GPU session \\
Session wall cap & default 30 minutes, maximum 1{,}440 & GPU session \\
Compute charge & GPU count $\times$ hours held $\times$ hardware multiplier,
in H100-equivalent hours & budget meter \\
Hardware multiplier & 1.0 for the H100 reference unit and for the GH200 nodes
we allocate on & budget meter \\
Wall warning and re-acquire threshold & 120\,s left on the allocation & GPU
session \\
Output returned from a step killed at the wall & last 4{,}000 characters & GPU
session \\
Grace before an ungranted allocation request is abandoned & the hold's wall
cap plus 8 hours, overridable from the environment & GPU session \\
Cores a CPU step is pinned to & 12 & sandbox \\
Address space of a CPU step & 16\,GiB & sandbox \\
Excerpt kept in a malformed report & 4{,}000 characters & report writer \\
\bottomrule
\end{tabular}
\caption{Harness parameters, their code defaults, and the component that
enforces each.}
\label{tab:knobs}
\end{table}

\subsection{Tools}
\label{app:toolset}

The model acts through the seven tools of Table~\ref{tab:tools}, and because
there is no file-reading tool and no web search, it reads files through the
shell with \texttt{grep}, \texttt{sed}, and \texttt{cat} and can fetch only
URLs it already knows or can construct.

\begin{table}[ht]
\centering
\small
\renewcommand{\arraystretch}{1.15}
\begin{tabular}{@{}l>{\raggedright\arraybackslash}p{9.4cm}@{}}
\toprule
Tool & Function \\
\midrule
\texttt{workspace\_bash} & Bash with the episode workspace as the working
directory, on the node that runs the harness, with no GPU passed through.
Cloning, installs, edits, inspection. Timeout default 900\,s, maximum
3{,}600\,s. Returns stdout, stderr, exit code, and duration. \\
\texttt{write\_file} & Write a UTF-8 file under the workspace or evidence roots,
up to 200{,}000 characters. \\
\texttt{edit\_file} & Replace one exact string in an existing file. The match
must be unique unless \texttt{replace\_all} is set. A bare repo-relative path is
located by suffix search when it resolves to exactly one file. \\
\texttt{update\_plan} & Replace the agent's ordered checklist of steps with
statuses \texttt{pending}, \texttt{in\_progress}, or \texttt{completed}. At most
one step may be \texttt{in\_progress}. \\
\texttt{fetch\_url} & Fetch a public http(s) URL and reduce HTML to text,
default 24{,}000 characters. \\
\texttt{list\_partitions} & Run \texttt{sinfo} and return each Slurm partition's
node counts by state, walltime limit, and GPU resources, plus the profile
default. Read-only. \\
\texttt{run\_gpu} & Run one command on the held GPU allocation
(Appendix~\ref{app:gpu-metering}). The only path to a GPU. \\
\bottomrule
\end{tabular}
\caption{The seven tools available to the model, with the code defaults and
caps each enforces.}
\label{tab:tools}
\end{table}

\subsection{Hardware setup and compute accounting}
\label{app:gpu-metering}

\texttt{run\_gpu} holds one Slurm allocation across calls, as
Figure~\ref{lst:slurm-lifecycle} shows. The first call acquires it with a fixed
GPU count and wall cap, later calls reuse it with no further queue wait, and
near the wall cap the harness releases and re-acquires. Queue time is not
charged, held time is charged including idle time while the model reasons, and
the harness refuses any allocation whose worst-case cost exceeds the remaining
budget. Charges apply between rounds for the wall time held since the previous
charge, and once the budget is exhausted the harness releases the session and
issues the final report turn. Step output streams to a numbered log under
\texttt{evidence/}, so a step killed at the wall still leaves output.

\begin{codebox}[lst:slurm-lifecycle]{bash}{Allocation lifecycle}
# acquire once; returns the instant the allocation is granted, then stays up
$ salloc --no-shell -A <account> -p <partition> --nodes=1 \
      --gpus=<k> --mem=<100*k>G --cpus-per-gpu=60 --time=<minutes>
# run each step into the held jobid, with no new queue wait
$ srun --jobid=<jobid> --ntasks=1 --unbuffered \
      apptainer exec --nv ... <image> bash -lc 'cd /repro/workspace && <cmd>'
# release when the agent is done, or at episode teardown
$ scancel <jobid>
\end{codebox}

\subsection{Environment details}
\label{app:sandbox}

Every shell step runs inside an Apptainer
container~\citep{kurtzer2017singularity} whose image is a read-only CUDA 12.9
and cuDNN development root with git, a C and C++ build chain, \texttt{uv}, and
Python 3.12 but no preinstalled PyTorch. Only bound paths are visible, the
workspace and evidence directories read-write at \texttt{/repro/workspace} and
\texttt{/repro/evidence}, the paper's reference copy read-only at
\texttt{/repro/reference}, and node-local \texttt{/tmp} and the shared package
caches read-write. The host home is not mounted, Apptainer runs with
\texttt{-{}-cleanenv} and \texttt{-{}-no-home}, and the harness forwards the
Hugging Face token, cache locations, and proxy settings through an environment
allowlist, so no secret appears on a command line.

\subsection{Run bundle and report format}
\label{app:report-format}

Each run writes a self-contained run bundle at
\texttt{<runs-dir>/\allowbreak<arxiv\_id>/\allowbreak<budget>h/\allowbreak<run\_id>/}.
The writable \texttt{workspace/} holds the agent's clones and scripts. The
read-only \texttt{reference/} holds the paper's LaTeX sources, its OpenReview
supplement, a file manifest, and an \texttt{info.json} of IDs. The
\texttt{evidence/} directory holds four fixed entries (\texttt{commands.log},
\texttt{trajectory.jsonl}, \texttt{env.lock}, and \texttt{patches/}), plus
one numbered log per GPU step, \texttt{plan.md}, and the \texttt{REPORT.md}
the task prompt requires. Per-run token counts and GPU-utilization samples are
recorded alongside, and the harness never re-executes the agent's scoring
command.

On the final, tool-free turn the model's output is constrained to a JSON
schema and saved as \texttt{report.json} beside \texttt{evidence/}, with the
fields of Table~\ref{tab:report-schema}. The harness overwrites
\texttt{paper\_id} with the ID it dispatched, so the model's echoed ID cannot
alter provenance, and records a
\texttt{report\_status} with the loop's
\texttt{exit\_reason}, and output that fails to parse or fails the structural
check is still written as a malformed report whose \texttt{report\_status} is
\texttt{degraded} and which includes the validation errors and an excerpt of
the raw output.

\begin{table}[ht]
\centering
\small
\renewcommand{\arraystretch}{1.15}
\begin{tabular}{@{}l>{\raggedright\arraybackslash}p{9.2cm}@{}}
\toprule
Field & Content \\
\midrule
\texttt{paper\_id} & The paper's arXiv ID, overwritten by the harness. \\
\texttt{claim} & The claim the agent worked toward, in the agent's own words. \\
\texttt{what\_ran} & What was executed, from the environment build to
the run. \\
\texttt{scoring\_command} & One command that reproduces the measured number from
a clean state. Empty when nothing ran. \\
\texttt{measurements} & Array of measured metrics. Each entry records the metric
name with units, the observed value, the paper's reference value or
\texttt{null}, the scope the value was measured over, and an array of citations
into \texttt{evidence/} supporting the observed value. \\
\texttt{agent\_assessment} & One of \texttt{reproduced}, \texttt{partial},
\texttt{not\_reproduced}, \texttt{could\_not\_run}. \\
\texttt{changes\_made} & Deviations from the released reference. \\
\texttt{blockers} & What blocked a reproduction, empty when nothing did. \\
\texttt{evidence\_files} & The files under \texttt{evidence/} the auditor should
open first. \\
\bottomrule
\end{tabular}
\caption{The nine fields of \texttt{report.json}, all required by the schema.}
\label{tab:report-schema}
\end{table}

Once the episode finishes the harness prunes the agent's virtual environment
and any oversized files from the workspace, and the auditor grades from
\texttt{report.json} and \texttt{evidence/} (Section~\ref{sec:auditing}).

\subsection{Prompt templates}
\label{app:prompts}

Figure~\ref{lst:harness-messages} in Appendix~\ref{app:prompt-listings} prints
every fixed string the model reads: the system message stored at index 0,
where compaction never reaches it, the final-turn instruction appended once
the harness removes the tools, and the five strings inserted during an
episode. The task prompt is the first user turn, a 290-line template whose
fourteen placeholder fields under nine headings are filled from the paper's
dataset row (Appendix~\ref{app:dataset-schema}), from the episode's compute
ceiling, or from the container paths of Appendix~\ref{app:sandbox}. Rendering
fails on an unfilled field, every field has a fallback string that states
absence, and the renderer omits \texttt{match\_target.config}, so the agent
derives the configuration that produces the target value from the paper.
Figure~\ref{lst:task-prompt} prints the template.

